\section{Where is the model fragile? A saliency-guided word attack}
\label{sec:saliency}

The word attack of \S\ref{sec:results} draws its substitution uniformly at
random from a small synonym table. It answers ``does \emph{any} word-level edit
break the model'' and leaves the distribution of fragility unmeasured: a flip
caused by deleting the word \emph{urgent} and a flip caused by replacing the word
\emph{the} are the same event in that accounting. This section asks the sharper
question, and pays for the answer in queries.

\paragraph{Method.} We implement the ordering of PWWS~\cite{ren2019pwws}. For a message
$x$ that the model reads as malicious with probability $p_0$, a content word $i$
and a candidate replacement $w$, let $p_{i,w} = P(\text{malicious} \mid x
\text{ with word } i \text{ replaced by } w)$ be measured with the same prompt.
Word saliency and the substitution priority are
\begin{equation}
S_i = p_0 - p_{i,w}, \qquad
\phi_{i,w} = \alpha \, \tilde{S}_i + (1-\alpha)\, p_{i,w},
\end{equation}
where $\tilde{S}_i$ is $S_i$ divided by the largest saliency in the message and
$\alpha = 0.5$. Substitutions are applied greedily in descending $\phi$; after
each accepted edit the verdict is re-queried, and the attack stops at the first
flip or when the shared 15\% word budget is exhausted.

The normalisation of $S$ is not cosmetic. With raw saliency and $\alpha = 0.5$,
$\phi_{i,w} = \tfrac{1}{2}(p_0 - p_{i,w}) + \tfrac{1}{2} p_{i,w} = \tfrac{1}{2}
p_0$ for \emph{every} candidate: the ranking collapses to a constant, and the
greedy loop then applies substitutions in the order the words happen to appear in
the message. Because $p_0$ is common to all candidates, only the scale of $S$
matters, and the scale is exactly what the normalisation fixes. Our first
implementation of this module had that defect; the unit test that now guards it
(\texttt{tests/check\_saliency.py}) asserts that a high-saliency word outranks a
neutral one at $\alpha = 0.5$.

\paragraph{Candidate replacements.} Two sources, recorded per substitution: a
curated security lexicon (66 entries: \emph{verify} $\rightarrow$
\emph{confirm/validate}, \emph{prize} $\rightarrow$ \emph{award/reward}, \ldots)
and WordNet~\cite{miller1995wordnet} synonyms ordered by lemma frequency. Frequency
ordering matters more than it looks: WordNet's third sense of a word is where
synonyms stop being paraphrases (\emph{tone} $\rightarrow$ \emph{timber}), and a
substitution like that would be counted as semantics-preserving when it is not.
Multi-word replacements are excluded so that the 15\% budget counts words and
nothing else. Candidates are capped at two per position and twelve positions per
message, which is what keeps the query cost bounded.

\paragraph{Two verdict channels.} This attack needs a number, not a paragraph: at
every candidate it asks the endpoint for $P(\text{malicious})$, obtained by
scoring the two permitted answer tokens in one forward pass. The other two
attacks read the model's generated one-word reply instead. Those are two
readouts of the same model, and comparing their ASRs without checking that they
agree would be comparing a ruler with a scale. Every stored saliency row is
therefore re-scored through the generation route, and Table~\ref{tab:saliency}
reports the generation verdict; the artefact keeps both columns and the
probability-route count, and we do not treat a disagreement between them as a
bug to be tuned away --- it is a fact about the model.

The probability route also has a failure mode worth recording, because it is
silent. The answer tokens are chosen per tokenizer, and Phi-3.5 encodes
``\,malicious'' as three tokens whose first element is the bare whitespace
marker that ``\,benign'' also begins with. Scoring the first token of each answer
therefore compares the space against itself: both logits are identical,
$P(\text{malicious}) = 0.5$ for every message, and a saliency attack built on it
finds nothing to rank and no flip to report. The endpoint now skips the shared
prefix, and raises rather than returning a constant when the two answers cannot
be separated. The saliency numbers for that model in this version come from the
corrected scoring; its earlier all-zero result was an artefact.

\paragraph{Query accounting.} Unlike the random attack, this one costs a measured
number of model calls: roughly two per probed position, plus one verification per
accepted edit. Probabilities come from a dedicated endpoint that scores the two
permitted answer tokens in a single forward pass over the prompt, batched across
candidates; the per-sample count is stored with the result, so the comparison in
Table~\ref{tab:saliency} is between attacks whose cost is on the record. The
asymmetry is real but smaller than an earlier draft of this paragraph claimed.
The random attack is not one draw per sample: it re-queries every original and
spends one query per attempt, up to three, which is 3.9 queries per message and
8{,}047 queries for the 2{,}079 messages it attacks. The saliency attack spends
22.9 per message --- 22{,}964 queries for 1{,}001 messages, six times as many
calls for 1.6 times as many flips. The paper therefore does not claim equal
cost, and it no longer claims the random attack was capped at one draw.

\paragraph{What the comparison isolates, and what it does not.} Three things
change between Table~\ref{tab:saliency}'s two rows at once: the \emph{ordering}
(uniform draw versus saliency), the \emph{candidate set} (a 14-entry hand-written
table versus a security lexicon plus WordNet), and the \emph{oracle} (a one-bit
verdict versus $P(\text{malicious})$). The table therefore supports only the
combination: saliency-ordered substitution over a security vocabulary breaks more
messages per unit of edit than a uniform draw over a small table. It does not
show that ordering alone is responsible, and it cannot say what the probability
access is worth. \S\ref{sec:ablation} specifies the experiment that separates
the three.

\IfFileExists{tables/t7_saliency.tex}{% Generated by src/make_paper_tex.py from results/attack_table.csv.
At the same word budget, the random attack flips 28 of 1035 caught messages (2.7\%) and the saliency attack flips 46 of 1001 (4.6\%).
Each saliency-attacked message costs a mean of 23 model queries. The random word attack is not a single draw either: it re-queries every original and spends one query per attempt, so its 1{,}035 messages cost 3{,}927 queries, 3.8 each. The saliency attacker buys its extra flips with about 6.1 times as many calls, and that price is part of the claim, not a footnote to it.
The saliency attack also edits MORE, not less: 9.3\% of words against 7.6\% for the random attack. Normalising by the budget actually spent, it buys 0.493 of a flip per percent of words edited against 0.355, so the advantage is real but smaller than the raw ASR ratio.
Both numbers are pooled over different candidate vocabularies and, after the exclusions, over different message sets, so the denominator is 1{,}001 for the saliency attack and 1{,}035 for the random one; the paired comparison of Section~\ref{sec:pairs} is the only place where the two are measured on the same messages.
Pooled query accounting, kept separate from the message count because the two are different numbers: the character attack sends 1{,}044 caught messages in 4{,}120 queries, the word attack 1{,}035 in 3{,}927, so the random attacks together send 2{,}079 messages in 8{,}047 queries (3.9 each), and the saliency attack sends 1{,}001 messages in 22{,}964 queries (22.9 each).
}{}

\begin{table*}[t]
\centering
\small
\caption{Saliency-guided (PWWS-style) word attack next to the uniform-random word attack, both inside the same 15\% word budget. Both spend a measured number of model queries: the random attack re-queries every original and keeps drawing until a flip or three attempts, the saliency attack probes every candidate. Per-sample query counts are stored with each row.}
\label{tab:saliency}
\begin{tabular}{llcccccc}
\toprule
\textbf{Task} & \textbf{Model} & \textbf{Attacked} & \textbf{Flips} & \textbf{ASR} & \textbf{95\% CI} & \textbf{Pert.} & \textbf{Jacc.} \\
\midrule
 SMS smishing & Phi-3.5-mini-instruct (saliency) & 163 & 12 & 7.4\% & 4.3--12.4 & 12.6\% & 0.773 \\
  & Phi-3.5-mini-instruct (word) & 161 & 6 & 3.7\% & 1.7--7.9 & 9.0\% & 0.918 \\
  & Qwen2.5-3B-Instruct (saliency) & 167 & 1 & 0.6\% & 0.1--3.3 & 12.3\% & 0.778 \\
  & Qwen2.5-3B-Instruct (word) & 186 & 6 & 3.2\% & 1.5--6.9 & 9.2\% & 0.916 \\
  & SmolLM2-1.7B-Instruct (saliency) & 185 & 11 & 5.9\% & 3.4--10.3 & 12.9\% & 0.768 \\
  & SmolLM2-1.7B-Instruct (word) & 184 & 3 & 1.6\% & 0.6--4.7 & 8.8\% & 0.916 \\
\midrule
 Phishing email & Phi-3.5-mini-instruct (saliency) & 120 & 14 & 11.7\% & 7.1--18.6 & 6.0\% & 0.867 \\
  & Phi-3.5-mini-instruct (word) & 118 & 6 & 5.1\% & 2.4--10.7 & 6.3\% & 0.945 \\
  & Qwen2.5-3B-Instruct (saliency) & 169 & 4 & 2.4\% & 0.9--5.9 & 5.4\% & 0.876 \\
  & Qwen2.5-3B-Instruct (word) & 190 & 5 & 2.6\% & 1.1--6.0 & 6.2\% & 0.950 \\
  & SmolLM2-1.7B-Instruct (saliency) & 197 & 4 & 2.0\% & 0.8--5.1 & 6.7\% & 0.851 \\
  & SmolLM2-1.7B-Instruct (word) & 196 & 2 & 1.0\% & 0.3--3.6 & 6.3\% & 0.948 \\
\bottomrule
\end{tabular}
\end{table*}


\subsection{A matched-vocabulary, matched-budget ablation}
\label{sec:ablation}

Table~\ref{tab:saliency} cannot attribute its difference to any single cause,
because ordering, candidate set and oracle all change together. The ablation
below holds two of the three fixed in every comparison, so each contrast has one
explanation. It is implemented in \texttt{src/ablation\_common.py}, exercised by
\texttt{tests/check\_ablation.py}, and run; its contrasts are
Table~\ref{tab:ablation} and the paragraph that discusses it, both generated from
\texttt{results/ablation\_contrasts.csv}. It is an ablation on a matched
subsample, not on the full attack set, and every count it reports is paired on
the messages the two arms share.

\paragraph{One probe phase, then a difference in policy.} Every arm attacks the
same messages, with the same candidate set, and pays the same probe cost before
it edits anything: one query for the original, then one query per single
substitution candidate, capped identically at the same number of calls per
message for every arm (24 in the configuration the runner reports, and the
cap it used travels with every row). The reported cells are the first 400
eligible messages of each model, and the cap never bound: the largest eligible
set is 382 messages (smollm2-1.7b), so every arm attacks exactly the messages the
attack tables report on, and an arm differs from another only in policy. The probe results are retained and only the \emph{policy} differs --- what
the arm does with those measurements. Matching the budget by construction rather
than by assumption is the point: a random-order arm pays exactly the 1 + $|C|$
calls a saliency arm pays before either starts substituting, so a difference in
flips cannot be a difference in cost. The realised count is still stored per row,
because a cap that is matched on paper and exhausted early in practice would be
the same error as an unmatched denominator.

\paragraph{Four cells and one stronger attack.} The factorial crosses the two
orderings with the two oracles:
\emph{random-prob} (uniform order, probability oracle), \emph{saliency-prob}
(the attack of this section, so the cell doubles as its re-implementation),
\emph{random-verdict} (uniform order, one-bit oracle) and \emph{group-verdict}
(a binary search over candidate \emph{subsets}: apply half, query once, recurse
into the half that flipped). A fifth arm, \emph{textfooler-prob}, implements
TextFooler's search~\cite{jin2020textfooler} rather than PWWS's: word importance
is leave-one-out --- each word replaced by a placeholder, importance the drop in
$P(\text{malicious})$, doubled when the leave-one-out text has already flipped,
which is the correction in the paper's Eq.~2 without which a word the model has
already lost looks unimportant --- then words are substituted in descending
importance with all of one word's candidates predicted in one batch, accepting
a flip that passes the meaning check and otherwise keeping the candidate that
lowers the malicious confidence most. It costs one extra query per position, so
it is matched by the cap rather than by construction and is reported as a
stronger reference point, not as a fifth cell of the factorial.

The contrasts are then read as: \emph{random-prob} vs \emph{saliency-prob} is
the ordering effect with vocabulary and budget held fixed; \emph{random-prob} vs
\emph{random-verdict} is what probability access is worth to an attacker who
cannot rank; \emph{random-verdict} vs \emph{group-verdict} is how much of that a
one-bit oracle can recover by searching instead of ranking; and
\emph{random-prob} against the recorded random word attack of
\S\ref{sec:results} isolates the candidate set, which is the only thing those
two do not share.

\paragraph{Meaning preservation is a constraint, not a claim.} Every arm gates
candidates through a lexical paraphrase test (WordNet-based, since this study
ships no sentence-embedding model) and records, per applied perturbation, the
similarity the gate used along with literal retention and paraphrase coverage.
TextFooler's arm turns the gate on by default at a recorded threshold, because
without a meaning constraint it is not TextFooler. All three scores travel with
every row, so a reader can re-derive an acceptance decision at another threshold
instead of taking ours.

\paragraph{What a one-bit attacker can and cannot do.} PWWS's ordering needs a
scalar per candidate. A verdict has one bit, and that bit is spent by the first
flip, so a verdict-only arm \emph{cannot rank} by saliency at all. Its targeted
policy is therefore a search over subsets, and the oracle contrast is reported as
how much of the saliency advantage survives when the attacker holds only the bit
--- not as a like-for-like ranking comparison, which would be measuring an
impossible attacker.

\paragraph{Calibration, seeds and status.} The meaning threshold and $\alpha$ are
tuned on the calibration slice, which is disjoint from both the evaluation and
the attack slice (\texttt{src/build\_splits.py} asserts the disjointness), and
never on the messages whose ASR is reported. Every arm takes a seed and the
sweep driver re-runs it at several, so ordering effects can be separated from
seed noise. One command per arm runs it
(\texttt{python src/run\_ablation.py --model <key> --arm <arm>}), and the
per-sample records are written in the same schema as the other attacks, which is
what lets the paired test of \S\ref{sec:pairs} be applied to the ablation
without new code.

% The table and its prose appear only when results/ablation_contrasts.csv
% exists, so the PDF cannot advertise cells no run produced.
\IfFileExists{tables/t9_ablation.tex}{%
% Generated by src/make_paper_tex.py from results/ablation_contrasts.csv.
Table~\ref{tab:ablation} reports each contrast on the messages its two arms share.
With the candidate set, the probe phase and the query cap held fixed, ordering by saliency flips 17 messages the uniform draw does not, against 0 in the other direction (exact McNemar $p=0.0000$ over 1001 paired messages).
Replacing the probability channel with the one-bit verdict loses 26 messages that survive on probabilities and wins 14 the probability arm misses ($p=0.0807$, 1001 pairs).
Searching over candidate subsets instead of ranking under the same bit recovers 6 messages the ranking arm misses, against 0 in the other direction ($p=0.0312$, 1044 pairs).
Against the 14-entry hand-written table, the lexicon plus WordNet flips 32 messages the small table does not, against 1 ($p=0.0000$, 1044 pairs).
TextFooler's leave-one-out importance ranking flips 0 messages PWWS's saliency ranking does not, against 20 in the other direction ($p=0.0000$, 1001 pairs). The published attack's ranking therefore is not uniformly stronger on these messages, which is why it is reported as a reference point rather than as a fifth cell.
These are pooled counts over model and task, and the cells reuse messages across models, so they describe a direction rather than a confirmatory test; the per-cell pairs are in the artefact.


\begin{table*}[t]
\centering
\small
\caption{Matched-vocabulary, matched-budget ablation of Section~\ref{sec:ablation}. Every contrast holds the candidate set, the probe phase and the 24-call cap fixed and varies one thing; A and B are the two arms, and positive A-only means the first arm flipped messages the second did not. Counts are pooled over model and task; the per-cell pairs are in \texttt{results/ablation\_contrasts.csv}.}
\label{tab:ablation}
\begin{tabular}{llcccccc}
\toprule
\textbf{Contrast} & \textbf{Isolates} & \textbf{Paired} & \textbf{A flips} & \textbf{B flips} & \textbf{A-only} & \textbf{B-only} & \textbf{$p$} \\
\midrule
 candidates & candidate set & 1044 & 45 & 14 & 32 & 1 & 0.0000 \\
 oracle & oracle & 1001 & 44 & 32 & 26 & 14 & 0.0807 \\
 ordering & ordering & 1001 & 61 & 44 & 17 & 0 & 0.0000 \\
 search & search under one bit & 1044 & 51 & 45 & 6 & 0 & 0.0312 \\
 stronger & search rule & 1001 & 41 & 61 & 0 & 20 & 0.0000 \\
\bottomrule
\end{tabular}
\end{table*}
}{}

\paragraph{Reading the five contrasts.} The ordering contrast is the one this
section was built for, and it is the cleanest result in the paper: with the
candidate set, the probe phase and the cap held fixed, every discordant message
went to saliency ordering. The candidate-set contrast is nearly as large again,
which is the uncomfortable half of the finding --- most of the difference between
Table~\ref{tab:saliency}'s two rows is vocabulary, not intelligence, so the
unmatched comparison overstates what ordering buys. Probability access, the
thing a saliency attack is supposed to need, is the weakest of the three: an
attacker holding only the one-bit verdict loses messages in one direction and
wins a few in the other, and the difference does not reach significance at this
sample size. Searching subsets instead of ranking recovers a handful of the
messages that bit costs, which is what the theory predicts, because one bit
cannot rank.

The fifth contrast needs its budget read out loud before its number is quoted.
TextFooler's leave-one-out importance spends 23.0 of the 24 calls per message
measuring word importance, where the saliency arm spends 20.6 and still has
edits left to spend its calls on; the extra query per position is what the
published search costs, and under a cap that has to cover both the probe and the
substitutions, it is paid for out of the substitutions. So the contrast says what
a query-starved importance ranking achieves inside this cap --- not that the
published attack is weaker, which this design cannot test. It is reported as a
reference point and not as a fifth cell for that reason, and the arm's
per-sample records carry the realised query count so a reader can check the
accounting rather than take it on faith.
